\scanpage{52}
\begin{flushright}
atkin \mspage{12}
\end{flushright}

\noindent which we now discuss \unsure{briefly|bríefly} in \handwritten{\S}\,6 below indicates that (5.4) is the right kind of property to generalise for all \uline{$q$}.

6. \uline{The} \uline{analogy} \uline{with} \uline{Hecke} \uline{theory}. We now return to the exact multiplicative properties such as (3.5). These have been extensively generalised by Hecke,Petersson,and others; for a useful summary we refer to RANKIN(196\hspace{.6em}). As a particular example of the theory we take the case of
\[
\handwritten{\sum_{n=2}^{\infty}}t_{2}(n)x^{n}=x^{2}f^{48}(x),
\]
for which Hecke showed that
\[
t_{2}(n)=a\,\handwritten{\lambda}(n)+\mathord{\rlap{$\handwritten{\bar{\phantom{a}}}$}a}\,\handwritten{\bar{\lambda}}(n),
\]
where $a,\handwritten{\lambda}(n)$,and their conjugates $\mathord{\rlap{$\handwritten{\bar{\phantom{a}}}$}a},\handwritten{\bar{\lambda}}(n)$ lie in the quadratic field $\deleted{\unsure{a|q}}$\reversemarginpar\margin{\handwritten{\tikz[x=1em,y=1em,baseline=-.1em]{\draw (.56,.7)--(.03,-.53) (.78,.39)--(.20,-.74); \draw (.33,.39) .. controls (.13,.46) and (.14,.16) .. (.27,.04) .. controls (.4,-.11) and (.26,-.26) .. (.25,-.3);}}} $Q(\handwritten{\surd}\,144169)$,and for $n\handwritten{\geqslant}1$ and prime $p$ we have
\[
\handwritten{\lambda}(np)-\handwritten{\lambda}(n)\handwritten{\lambda}(p)+p^{23}\handwritten{\lambda}(n/p)=0.
\tag*{(6.1)}
\]

It seems clear that the multiplicative congruence properties we have described in terms of extensions of residue class rings are closely anal\deleted{g}\reversemarginpar\margin{\inserted{\handwritten{\unsure{a|u}}} \handwritten{\tikz[x=1em,y=1em,baseline=-.1em]{\draw (.10,-.45)--(.47,1.10);}}}gous to the exact multiplicative properties of Hecke theory in terms of algebraic number fields. We have not described in this account the underlying common feature of the Hecke operators,but it seems reasonable to conjecture that some kind of ``$\handwritten{\mathfrak{p}}$-adic'' Hecke theory\normalmarginpar\margin{\handwritten{$\left(\begin{gathered}\text{German}\\[-.4ex]p\end{gathered}\right)$}} exists for modular forms of positive and(as for $c(n)$) zero dimension, which is a counterpart of the exact theory for forms of negative dimension. The present difficulties in the way of precise formulation and proof see\reversemarginpar\margin{\inserted{\handwritten{m}} \handwritten{\tikz[x=1em,y=1em,baseline=-.1em]{\draw (.08,-.55)--(.57,1.05); \draw (.19,-.24) .. controls (.30,-.18) and (.29,-.49) .. (.58,-.50);}}}\rlap{\smash{\handwritten{\tikz[x=1em,y=1em,baseline=-.1em]{\draw (.08,-.55)--(.57,1.05); \draw (.19,-.24) .. controls (.30,-.18) and (.29,-.49) .. (.58,-.50);}}}} formidable in so far as the \uline{congruences} we have considered in this paper arise essentially at just one point\normalmarginpar\margin{\handwritten{\tikz[x=1em,y=1em,baseline=-.1em]{\draw (.08,-.55)--(.57,1.05); \draw (.19,-.24) .. controls (.30,-.18) and (.29,-.49) .. (.58,-.50);}} \inserted{\handwritten{$\tau=$}}}($\rlap{\smash{\handwritten{\tikz[x=1em,y=1em,baseline=-.1em]{\draw (.08,-.55)--(.57,1.05); \draw (.19,-.24) .. controls (.30,-.18) and (.29,-.49) .. (.58,-.50);}}}}i\handwritten{\infty}$) of the complex upper half-plane,whereas the \uline{identities} of Hecke theory exist at all points of the upper half-plane,taking the form of multiplicative properties for Fourier coefficients at\normalmarginpar\margin{\handwritten{\tikz[x=1em,y=1em,baseline=-.1em]{\draw (.08,-.55)--(.57,1.05); \draw (.19,-.24) .. controls (.30,-.18) and (.29,-.49) .. (.58,-.50);}} \inserted{\handwritten{$\tau=$}}} $\rlap{\smash{\handwritten{\tikz[x=1em,y=1em,baseline=-.1em]{\draw (.08,-.55)--(.57,1.05); \draw (.19,-.24) .. controls (.30,-.18) and (.29,-.49) .. (.58,-.50);}}}}i\handwritten{\infty}$.

7. \uline{Non-Ramanujan} \uline{congruences} \uline{for} \uline{primes} $q\handwritten{\leqslant}11$. Once the above properties for primes $q\handwritten{\geqslant}13$ had been found it was clear that something similar must
% archive sheet number: 51
